\section{系统设计}\label{ux7cfbux7edfux8bbeux8ba1}

本章在需求分析的基础上说明系统总体设计。前一章已经将课题目标细化为任务接入、候选计划生成、工具执行、HITL、运行时恢复和审计追踪等需求。基于这些需求，本文系统设计的重点不是把若干工具串接成固定流水线，而是在可执行约束、安全约束和预算约束下，为蛋白质设计任务提供可规划、可恢复、可审计的工作流运行框架。

本章先介绍系统分层架构、核心组件、有限状态机和六阶段蛋白质设计工作流，再定义约束与证据感知、信念引导、恢复自适应工作流规划（Constraint-
and Evidence-aware Belief-guided Recovery-adaptive Workflow
Planning，CEBRA-WP）。针对
CEBRA-WP，本章给出其在本文系统中的形式化对象、评分函数、运行时状态、恢复动作和实验可验证性。

\subsection{设计目标与总体架构}\label{ux8bbeux8ba1ux76eeux6807ux4e0eux603bux4f53ux67b6ux6784}

蛋白质设计工作流与一般 Web
后端任务存在明显差异。任务目标具有科学语义，用户可能只给出``设计一个稳定短肽''或``提高候选序列结构质量''等高层目标，系统需要将其转换为可执行步骤。工具调用的成本差异也较大：序列质量检查通常较轻，结构预测、重型打分和重设计则可能带来更高计算成本。与此同时，中间失败不一定意味着整个任务失败，参数错误、输入引用缺失、质量门禁未通过和外部工具暂时不可用，都可能通过局部修补或后缀重规划恢复。系统输出还需要面向科研场景保留证据链，包括候选生成依据、人工确认记录、执行结果、失败原因和恢复历史。

基于上述特点，系统采用五层分层架构，如图 4-1 所示。

图 4-1 系统五层分层架构
插图文件：\texttt{paper/figures/system-architecture.drawio.svg}

输入层面向 Web 工作台、CLI 和
API，负责接收自然语言目标、结构化约束和人工决策。智能规划层以
PlannerAgent、ProteinToolKG 和 CEBRA-WP
策略为核心，负责生成候选计划、进行可行性过滤、评分和恢复候选生成。执行层由
ExecutorAgent、PlanRunner 和 StepRunner
组成，负责按已确认计划调用工具并记录步骤结果。安全与汇总层包括
SafetyAgent 和
SummarizerAgent，分别负责风险判定与最终报告生成。资源层包括 ToolAdapter
注册表、ProteinToolKG、事件日志、任务快照和文件产物管理。

五层之间通过结构化数据契约交换信息。输入层不直接执行工具，规划层不直接修改任务终态，执行层不越过有限状态机进行状态跳转，安全层不直接编辑计划，汇总层不重新执行计算。通过这种职责拆分，候选生成、工具调用、人工确认和恢复动作都能回到对应的契约对象与事件记录，从而保证复杂任务中的可追溯性。

\subsection{核心组件与职责边界}\label{ux6838ux5fc3ux7ec4ux4ef6ux4e0eux804cux8d23ux8fb9ux754c}

系统核心逻辑由四类 Agent 和工具适配层共同完成。

PlannerAgent
负责计划生成和恢复候选生成。它读取任务目标、约束、执行历史和
ProteinToolKG，生成初始
PlanCandidate，或在失败后生成局部修补候选与后缀重规划候选。PlannerAgent
的输出必须是结构化候选，而不是无法验证的自然语言建议。每个候选包含候选标识、摘要、结构化载荷、评分分解、风险等级、成本估计、解释文本和来源引用，因此既能被自动策略排序，也能呈现在人工审查界面中。

ExecutorAgent 是工具调用与计划推进组件。它通过 PlanRunner
管理计划级执行流程，通过 StepRunner
处理单个步骤的输入解析、上游引用求解、适配器调用和结果写入。ExecutorAgent
可以识别失败、重试耗尽和安全阻断等信号，但不自行绕过 Planner
生成修补方案，也不在等待人工确认时继续执行工具。

SafetyAgent 是风险信号源。它在输入、执行过程和输出阶段产生安全判定，输出
ok、warn、block 等等级。warn 可触发人工确认，block
可阻断自动推进并触发重规划候选。SafetyAgent
的定位是风险判定与建议，不负责计划搜索和工具执行。

SummarizerAgent
是结果汇总组件。它读取计划、步骤结果、安全事件和恢复历史，生成面向用户的报告和机器可读的
DesignResult。其输出应反映已完成计算和已有证据，不把未验证推断写成实验结论。

ToolAdapter
层为外部工具提供统一调用接口。工具注册表面向执行层，回答``如何调用工具''；ProteinToolKG
面向规划层，回答``哪些工具可用于何种能力、输入输出如何兼容、成本和风险如何估计''。二者共同支撑候选生成与执行验证。ProteinToolKG
的局部结构如图 4-2 所示。

图 4-2 ProteinToolKG 局部可视化
插图文件：\texttt{paper/figures/protein-toolkg-local-view.drawio.svg}

图 4-2 展示了 ProtGPT2、ProteinMPNN、ESMFold/OpenFold、Biopython
QC、DSSP 和 objective\_ranker
等代表性工具之间的能力节点、输入输出字段、成本风险属性和 I/O
兼容边。该图强调本文使用的 ProteinToolKG
是轻量工具能力索引，而非完整生物知识图谱。它的主要作用是为 Planner
提供能力发现、约束校验和候选解释依据。

\subsection{任务生命周期、FSM 与 HITL
约束}\label{ux4efbux52a1ux751fux547dux5468ux671ffsm-ux4e0e-hitl-ux7ea6ux675f}

系统使用有限状态机（Finite State Machine，FSM）控制任务生命周期，如图
4-3 所示。

图 4-3 FSM 状态转移图
插图文件：\texttt{paper/figures/fsm-state-transition.drawio.svg}

任务对外状态包括
CREATED、PLANNING、WAITING\_PLAN\_CONFIRM、PLANNED、RUNNING、WAITING\_PATCH\_CONFIRM、WAITING\_REPLAN\_CONFIRM、SUMMARIZING、DONE、FAILED
和 CANCELLED。WAITING\_*
状态具有明确语义：系统已经暂停自动推进，等待人类提交结构化
Decision。DONE、FAILED 和 CANCELLED 为终态，进入后不再被自动修改。

人在环决策（Human-in-the-loop，HITL）不是任意插入的人机交互按钮，而是由成本、风险和不确定性共同触发的受控状态。不同等待状态的进入条件见图
4-4。

图 4-4 HITL 决策条件
插图文件：\texttt{paper/figures/hitl-decision-conditions.drawio.png}

PLANNING 阶段，如果候选置信度不足、即将调用高代价工具或 SafetyAgent 给出
warn，系统进入 WAITING\_PLAN\_CONFIRM。RUNNING 阶段，如果步骤失败但仍有
retry budget，系统先进行有界重试；若重试耗尽且局部可修复，则进入
WAITING\_PATCH\_CONFIRM；若出现结构性失败、安全阻断或恢复余量不足，则进入
WAITING\_REPLAN\_CONFIRM。人工决策通过 Decision 契约提交，Decision Apply
模块验证 pending action 与候选绑定关系，再推动 FSM 合法迁移。

这一设计为后续 CEBRA-WP
算法提供了控制边界：算法可以生成候选、计算分数和提出动作建议，但任务状态仍由
FSM 统一推进；算法建议不能绕过 WAITING\_* 状态，也不能把 stop
直接变成不可审计的隐式失败。

\subsection{六阶段蛋白质设计工作流}\label{ux516dux9636ux6bb5ux86cbux767dux8d28ux8bbeux8ba1ux5de5ux4f5cux6d41}

蛋白质设计能力被组织为六个阶段，如图 4-5 所示。

图 4-5 六阶段 de novo 蛋白质设计工作流
插图文件：\texttt{paper/figures/workflow-flowchart.drawio.svg}

六个阶段分别为序列探索、结构映射、质量门禁、结构条件精修、目标评分和结果汇总。序列探索生成候选序列；结构映射将序列转换为结构并产生折叠置信度；质量门禁检查序列合法性、结构完整性和低复杂度等工程质量；结构条件精修根据结构反馈进行重设计；目标评分对候选进行多目标排序；结果汇总生成
DesignResult 和报告。

六阶段是能力分层，而不是不可改变的线性流水线。质量门禁失败可以回到序列探索，目标评分不足可以回到结构条件精修，结构映射与精修之间可以形成迭代闭环。CEBRA-WP
的作用正是在这些可替代路径之间进行约束化选择：当证据不足时延后高代价步骤，当局部失败可修复时优先
\texttt{patch\_local}（局部修补），当后缀不再可靠时保留有效前缀并进行
\texttt{suffix\_replan}（后缀重规划）。

\subsection{CEBRA-WP
的提出背景}\label{cebra-wp-ux7684ux63d0ux51faux80ccux666f}

上一节给出了蛋白质设计工作流的能力分层和可替代路径，但阶段划分本身并不能回答何时继续执行、何时局部修补、何时后缀重规划以及何时终止止损。CEBRA-WP
的提出，正是为了处理这些运行时决策。

传统固定流水线适合步骤稳定、成本均匀、失败语义简单的任务，但难以覆盖本文场景。蛋白质设计工作流中的关键风险往往具有部分可观测性：一次结构预测失败可能来自工具服务异常，也可能表明上游序列质量不足；一次质量门禁失败可能通过阈值、参数或局部工具替换修复，也可能说明当前候选链路整体失效。系统若只按静态计划执行，会在证据不足时过早调用高代价工具；若只依赖单次
LLM 规划，候选又缺少可复现的约束过滤和恢复依据。

CEBRA-WP
的设计动机主要来自三类研究。部分可观测规划研究指出，决策不应只依赖瞬时观测，而应维护对隐藏状态的
belief-state
表示\citep{kaelbling1998pomdp, shani2024heuristics}。在本文中，隐藏状态不是完整环境模型，而是当前链路是否仍有成功可能、失败是否具有结构性、剩余预算压力和证据是否充分等低维量。预算约束规划与预算强化学习研究强调，成本和风险应作为一等决策变量，而不是执行失败后的附加统计\citep{carrara2019budgetedrl}；这与蛋白质设计中高代价结构预测和目标评分的治理需求相一致。Tree
of Thoughts
强调保留多个候选推理路径并进行搜索\citep{yao2023tot}，Reflexion
则强调利用失败反馈改进后续任务表现\citep{shinn2023reflexion}。本文将这两类思想约束到可执行工具链规划中，使候选、恢复动作和人工确认都以结构化对象表达。

基于这些考虑，CEBRA-WP
被定义为一种面向高代价科研工作流的约束化、证据感知、信念引导、恢复自适应规划算法。它不追求训练端到端最优控制器，而是以可解释、可配置、可审计的形式，把候选生成、硬约束过滤、静态效用、运行时状态、后验证据和恢复动作选择组织为一个闭环。

\subsection{CEBRA-WP
形式化定义}\label{cebra-wp-ux5f62ux5f0fux5316ux5b9aux4e49}

CEBRA-WP 的完整名称为 Constraint- and Evidence-aware Belief-guided
Recovery-adaptive Workflow
Planning，本文译为``约束与证据感知、信念引导、恢复自适应的工作流规划''。当前论文版本为
\texttt{cebra\_wp.v2}，其子公式包括
\texttt{static\_score.v1}、\texttt{posterior\_score.v1}、\texttt{runtime\_adjustment.v1}、\texttt{action\_utility.v1}
和 \texttt{action\_bias.v1}。

算法总体闭环如图 4-6 所示。

图 4-6 CEBRA-WP 算法闭环
插图文件：\texttt{paper/figures/algorithm-loop.drawio.svg}

图 4-6 的上半部分是静态规划层：系统根据目标、约束、ProteinToolKG
和历史记录生成候选工具链，随后执行硬可行性过滤与静态评分。图 4-6
的下半部分是运行时自适应层：系统根据
StepResult、SafetyResult、失败上下文、预算消耗和恢复历史更新 Lite
belief-state / 轻量信念状态，再对候选进行运行时重排序，并在
\texttt{continue}、\texttt{patch\_local}、\texttt{suffix\_replan} 和
\texttt{stop} 之间选择动作建议。

算法使用的主要符号如表 4-1 所示。

\textbf{表 4-1 CEBRA-WP 符号与语义定义}

{\def\LTcaptype{none} % do not increment counter
\begin{longtable}[]{@{}
  >{\raggedright\arraybackslash}p{(\linewidth - 6\tabcolsep) * \real{0.2500}}
  >{\raggedright\arraybackslash}p{(\linewidth - 6\tabcolsep) * \real{0.2500}}
  >{\raggedright\arraybackslash}p{(\linewidth - 6\tabcolsep) * \real{0.2500}}
  >{\raggedright\arraybackslash}p{(\linewidth - 6\tabcolsep) * \real{0.2500}}@{}}
\toprule\noalign{}
\begin{minipage}[b]{\linewidth}\raggedright
符号
\end{minipage} & \begin{minipage}[b]{\linewidth}\raggedright
含义
\end{minipage} & \begin{minipage}[b]{\linewidth}\raggedright
来源或载体
\end{minipage} & \begin{minipage}[b]{\linewidth}\raggedright
在算法中的作用
\end{minipage} \\
\midrule\noalign{}
\endhead
\bottomrule\noalign{}
\endlastfoot
\texttt{g} & 设计目标 & 用户输入、ConfirmedTaskSpec &
定义任务目标与目标权重 \\
\texttt{C} & 约束集合 & 任务约束、策略配置 &
限定长度、安全、预算、工具白名单和输出要求 \\
\texttt{K} & ProteinToolKG & ProteinToolKG & 提供工具能力、I/O
schema、兼容关系、成本和风险 \\
\texttt{h\_t} & 时间步 \texttt{t} 前的执行历史 & EventLog、TaskSnapshot
& 提供已完成步骤、失败记录、恢复历史和人工决策 \\
\texttt{o\_t} & 当前运行时观测 &
StepResult、SafetyResult、指标、错误细节 &
更新运行时状态并计算后验证据 \\
\texttt{x\_t} & Lite belief-state / 轻量信念状态 & RuntimeState &
表示成功概率、结构性失败概率、恢复余量、剩余成本和证据充分度 \\
\texttt{Pi\_raw,t} & 原始候选集合 & PlannerAgent &
包含初始计划、局部修补或重规划候选 \\
\texttt{Pi\_t} & 过滤后的候选集合 & FeasibilityFilter &
只保留可执行或受保护的 degraded feasible 候选 \\
\texttt{S\_static} & 静态效用 & score\_breakdown &
在运行时观测介入前评价候选先验质量 \\
\texttt{G\_post} & 后验目标匹配 & posterior\_objective &
根据证据可靠性修正目标评分 \\
\texttt{Delta} & 运行时修正项 & runtime\_adjustment &
根据状态变量修正候选排序 \\
\texttt{U\_pi} & 候选运行时效用 & RuntimeEvaluator & 输出 Top-K
和默认候选 \\
\texttt{a\_t} & 恢复动作 & action\_utility & 在
\texttt{continue}、\texttt{patch\_local}、\texttt{suffix\_replan}、\texttt{stop}
之间选择 \\
\end{longtable}
}

在时间步 \texttt{t}，算法输入为
\texttt{(g,\ C,\ K,\ h\_t,\ o\_t,\ x\_t)}，输出为结构化 Decision
建议：候选集合 \texttt{Pi\_t}、默认候选
\texttt{pi*}、候选解释、运行时状态摘要和恢复动作
\texttt{a\_t}。核心计算过程可概括为：

\begin{Shaded}
\begin{Highlighting}[]
\NormalTok{Pi\_raw,t = GenerateCandidates(g, C, K, h\_t)}
\NormalTok{Pi\_t     = FeasibilityFilter(Pi\_raw,t, C, K, h\_t)}
\NormalTok{S\_static = StaticUtility(pi, g, C, K)}
\NormalTok{x\_t+1    = BeliefUpdate(x\_t, o\_t, h\_t)}
\NormalTok{G\_post   = PosteriorObjective(pi, g, o\_t)}
\NormalTok{U\_pi     = RuntimeCandidateUtility(S\_static, G\_post, x\_t+1)}
\NormalTok{a\_t      = RecoveryAwareActionSelection(x\_t+1, Pi\_t, h\_t, C)}
\end{Highlighting}
\end{Shaded}

该定义突出两点：一是候选生成与候选选择相互分离，Planner
可以生成多个方案，但只有通过可行性和效用评估的候选才会进入执行或人工确认；二是运行时状态只修正排序和动作建议，不改变
FSM、Agent 职责和硬约束。

\subsection{候选生成、硬可行性过滤与静态效用}\label{ux5019ux9009ux751fux6210ux786cux53efux884cux6027ux8fc7ux6ee4ux4e0eux9759ux6001ux6548ux7528}

CEBRA-WP 支持三类候选。PlanCandidate 表示初始完整计划；PatchCandidate
表示对当前计划中局部步骤的参数级、工具级或结构级修补；ReplanCandidate
表示对未执行后缀或整体策略的替换，其中 \texttt{suffix\_replan}
优先保留已验证前缀，\texttt{terminal\_stop}
表示继续投入不划算时的终止型重规划候选。

候选进入评分前必须先通过硬可行性过滤。对候选 \texttt{pi}
定义硬可行性谓词：

\begin{Shaded}
\begin{Highlighting}[]
\NormalTok{F\_h(pi, C, K, h\_t)}
\NormalTok{  = F\_tool ∧ F\_schema ∧ F\_io ∧ F\_safety ∧ F\_budget{-}hard ∧ F\_availability}
\end{Highlighting}
\end{Shaded}

其中，\texttt{F\_tool}
检查候选中的工具是否存在于能力图或适配器注册表；\texttt{F\_schema}
检查输入输出字段是否满足工具 schema；\texttt{F\_io}
检查跨步骤引用是否闭合；\texttt{F\_safety} 检查是否违反安全等级或触发
safety block；\texttt{F\_budget-hard}
检查是否突破不可逾越预算上限；\texttt{F\_availability}
检查关键工具是否可用或是否有明确降级路径。过滤后的集合为：

\begin{Shaded}
\begin{Highlighting}[]
\NormalTok{Pi\_t = \{ pi in Pi\_raw,t | F\_h(pi, C, K, h\_t) = 1 \}}
\end{Highlighting}
\end{Shaded}

硬可行性过滤承担系统边界约束，因此静态评分不能把硬不可行候选``打高分后救回''。工程实现可以保留
degraded feasible 候选用于解释或 HITL
审查，但这类候选必须携带降级原因和人工确认要求。

在候选通过硬过滤后，算法计算静态效用：

\begin{Shaded}
\begin{Highlighting}[]
\NormalTok{S\_static(pi)}
\NormalTok{  = w\_f F\_s(pi)}
\NormalTok{  + w\_g G(pi; g, o\_t)}
\NormalTok{  {-} w\_c C\_norm(pi)}
\NormalTok{  {-} w\_r R\_norm(pi)}
\NormalTok{  {-} w\_rec Rec(pi)}
\NormalTok{  + w\_q Q(pi)}
\end{Highlighting}
\end{Shaded}

\texttt{F\_s} 是软可行性分数，表示候选在 schema 完整、工具
readiness、fallback depth 等方面的先验质量；\texttt{G}
是目标匹配度；\texttt{C\_norm} 是归一化成本；\texttt{R\_norm}
是归一化风险；\texttt{Rec} 是恢复复杂度；\texttt{Q} 是工程可靠性项。权重
\texttt{w\_*}
由策略配置给出，用于在不同任务设置下调节目标、成本、风险和恢复难度的相对重要性。

静态效用用于形成执行前的候选先验排序。它可以偏好成本更低、风险更小、恢复更容易的工具链，但无法感知执行中已经发生的失败、证据不足或预算压力变化。因此，静态效用之后还需要
Lite belief-state / 轻量信念状态和后验证据修正。

\subsection{Lite belief-state
与观测更新}\label{lite-belief-state-ux4e0eux89c2ux6d4bux66f4ux65b0}

CEBRA-WP 使用 Lite belief-state /
轻量信念状态近似表示执行过程中无法完全观测的工作流状态。它只保留三类变量：对恢复动作选择必要的变量、能够从现有日志稳定更新的变量，以及能够被实验和案例解释的变量。状态向量定义为：

\begin{Shaded}
\begin{Highlighting}[]
\NormalTok{x\_t = [}
\NormalTok{  p\_success,}
\NormalTok{  p\_structural\_failure,}
\NormalTok{  recovery\_margin,}
\NormalTok{  expected\_remaining\_cost,}
\NormalTok{  evidence\_sufficiency}
\NormalTok{]}
\end{Highlighting}
\end{Shaded}

各状态量的语义如表 4-2 所示。

\textbf{表 4-2 Lite belief-state / 轻量信念状态变量}

{\def\LTcaptype{none} % do not increment counter
\begin{longtable}[]{@{}
  >{\raggedright\arraybackslash}p{(\linewidth - 8\tabcolsep) * \real{0.2000}}
  >{\raggedright\arraybackslash}p{(\linewidth - 8\tabcolsep) * \real{0.2000}}
  >{\raggedright\arraybackslash}p{(\linewidth - 8\tabcolsep) * \real{0.2000}}
  >{\raggedright\arraybackslash}p{(\linewidth - 8\tabcolsep) * \real{0.2000}}
  >{\raggedright\arraybackslash}p{(\linewidth - 8\tabcolsep) * \real{0.2000}}@{}}
\toprule\noalign{}
\begin{minipage}[b]{\linewidth}\raggedright
状态量
\end{minipage} & \begin{minipage}[b]{\linewidth}\raggedright
取值范围
\end{minipage} & \begin{minipage}[b]{\linewidth}\raggedright
含义
\end{minipage} & \begin{minipage}[b]{\linewidth}\raggedright
主要更新来源
\end{minipage} & \begin{minipage}[b]{\linewidth}\raggedright
决策作用
\end{minipage} \\
\midrule\noalign{}
\endhead
\bottomrule\noalign{}
\endlastfoot
\texttt{p\_success} & \texttt{{[}0,1{]}} &
当前链路继续执行后完成任务的估计概率 &
步骤成功、质量指标、候选静态分、失败记录 & 支持 \texttt{continue} 与
\texttt{stop} 判断 \\
\texttt{p\_structural\_failure} & \texttt{{[}0,1{]}} &
当前链路遭遇结构性失败或后续必须重规划的估计概率 &
结构预测失败、质量门禁失败、安全阻断、重复失败 & 提高
\texttt{patch\_local}、\texttt{suffix\_replan}、\texttt{stop} 权重 \\
\texttt{recovery\_margin} & \texttt{{[}0,1{]}} &
在保留有效前缀前提下继续恢复的余量 &
已完成步骤比例、失败类型、patch/replan 次数 &
区分局部修补与后缀重规划 \\
\texttt{expected\_remaining\_cost} & 非负实数 &
从当前状态到任务结束的剩余成本暴露 &
剩余步骤、工具成本先验、预算配置、重试记录 &
派生预算压力并约束高代价动作 \\
\texttt{evidence\_sufficiency} & \texttt{{[}0,1{]}} &
当前证据是否足以支持进入更高代价步骤 &
质量门禁、结构指标、目标评分、证据可靠性 &
控制高代价步骤推进与人工确认 \\
\end{longtable}
}

运行时观测 \texttt{o\_t} 来自
StepResult、SafetyResult、失败上下文、局部修补/重规划历史、已完成步骤、剩余后缀和
HITL 决策记录。派生量如
\texttt{budget\_pressure}、\texttt{intervention\_value}、\texttt{local\_patchability}
和 \texttt{prefix\_preservability}
不作为持久化主状态，而是根据当前状态和候选上下文按需计算。这样可以避免把与具体候选强相关的临时判断写入长期状态，从而降低状态漂移风险。

以预算压力为例，\texttt{expected\_remaining\_cost}
保留剩余成本原始估计，不直接等同于 \texttt{{[}0,1{]}}
区间内的预算压力。若任务给出预算上限 \texttt{budget\_cap}，则：

\begin{Shaded}
\begin{Highlighting}[]
\NormalTok{budget\_pressure}
\NormalTok{  = clip(expected\_remaining\_cost / max(budget\_cap, 0.1), 0, 1.5)}
\end{Highlighting}
\end{Shaded}

若任务未给出预算上限，则使用：

\begin{Shaded}
\begin{Highlighting}[]
\NormalTok{budget\_pressure = clip(expected\_remaining\_cost, 0, 1.5)}
\end{Highlighting}
\end{Shaded}

因此，预算压力是由任务上下文派生的动作决策量，而不是 RuntimeState
的持久化字段。该区分使系统既能记录可解释的剩余成本，也能在不同预算设置下统一比较恢复动作。

状态初始化可由首选候选的静态分数、风险分数、恢复复杂度、剩余成本先验和廉价证据覆盖率给出。例如，\texttt{p\_success}
可由静态候选分数裁剪得到，\texttt{p\_structural\_failure}
可由风险项裁剪得到，\texttt{recovery\_margin}
可由恢复复杂度的反向量估计，\texttt{expected\_remaining\_cost}
来自剩余工具成本先验，\texttt{evidence\_sufficiency}
来自低成本质量证据覆盖率。随着执行推进，成功步骤会提升
\texttt{p\_success} 和证据充分度，结构性失败会提高
\texttt{p\_structural\_failure}
并降低恢复余量，重复失败和预算消耗则会提高剩余成本压力。

该状态设计吸收了部分可观测规划中 belief-state
的思想\citep{kaelbling1998pomdp}，但在工程上保持轻量：系统不求解完整
POMDP，也不在线训练策略网络，而是使用可解释变量为候选重排和恢复动作提供一致依据。

\subsection{后验目标评分与运行时重排序}\label{ux540eux9a8cux76eeux6807ux8bc4ux5206ux4e0eux8fd0ux884cux65f6ux91cdux6392ux5e8f}

蛋白质设计目标通常是多目标的，包括结构质量、稳定性、功能、novelty
和安全性等。不同目标的证据来源不同，证据可靠性也不同。为避免把弱证据与直接证据等价处理，CEBRA-WP
使用证据加权后验目标匹配：

\begin{Shaded}
\begin{Highlighting}[]
\NormalTok{G\_post(pi; g, o\_t) = Σ\_m λ\_m(g) · ρ\_m(o\_t) · q\_m(pi, o\_t)}
\end{Highlighting}
\end{Shaded}

其中，\texttt{m} 表示目标维度，\texttt{λ\_m(g)}
是由任务目标决定的目标权重，\texttt{q\_m(pi,\ o\_t)}
是候选在该维度上的归一化分数，\texttt{ρ\_m(o\_t)}
是证据可靠性权重。证据状态分为 direct、proxy、degraded 和
missing：direct evidence 指已产生的结构质量指标，proxy evidence
指由轻量质量门禁推断出的间接信号，degraded evidence
表示工具降级或证据覆盖不完整，missing
表示当前没有可用证据。整体证据充分度可写为：

\begin{Shaded}
\begin{Highlighting}[]
\NormalTok{e\_t = clip(Σ\_m λ\_m(g) · ρ\_m(o\_t), 0, 1)}
\end{Highlighting}
\end{Shaded}

该值进入
\texttt{evidence\_sufficiency}，影响是否继续进入高代价步骤。若证据不足而候选下一步成本较高，运行时重排序会降低该候选；若已有充分低成本证据支持，则候选可以更合理地进入结构预测或目标评分。

将 \texttt{G\_post}
替换静态效用中的目标匹配项后，得到包含后验证据的基础分
\texttt{S\_post}。运行时重排序定义为：

\begin{Shaded}
\begin{Highlighting}[]
\NormalTok{U\_pi(pi, x\_t) = clip(S\_post(pi) + Delta(pi, x\_t), 0, 1)}
\end{Highlighting}
\end{Shaded}

其中 \texttt{Delta(pi,\ x\_t)} 是有界运行时修正项，取值范围控制在
\texttt{{[}-0.35,\ 0.35{]}}。其一般形式为：

\begin{Shaded}
\begin{Highlighting}[]
\NormalTok{Delta(pi, x\_t)}
\NormalTok{  = k\_s · (p\_success {-} 0.5) · Conf(pi)}
\NormalTok{  + k\_e · (2 · evidence\_sufficiency {-} 1) · max(Conf(pi), F\_s(pi))}
\NormalTok{  {-} k\_f · p\_structural\_failure · (1 {-} RiskScore(pi))}
\NormalTok{  + k\_r · recovery\_margin · RecoveryScore(pi)}
\NormalTok{  {-} k\_c · budget\_pressure · (1 {-} CostScore(pi))}
\NormalTok{  + k\_a · ActionBias(pi, x\_t)}
\end{Highlighting}
\end{Shaded}

其中，\texttt{Conf(pi)} 表示候选置信度，\texttt{RiskScore(pi)} 和
\texttt{CostScore(pi)}
为候选的风险与成本质量分数，\texttt{RecoveryScore(pi)}
表示候选恢复友好度，\texttt{ActionBias(pi,\ x\_t)}
表示候选与当前恢复动作偏好的匹配程度。其直观含义是：当
\texttt{p\_structural\_failure}
和预算压力升高时，算法降低高成本、低可恢复候选的排序；当
\texttt{p\_success}、\texttt{recovery\_margin} 和
\texttt{evidence\_sufficiency}
较高时，算法提高继续执行或低风险候选的排序。

运行时重排序有两个边界：一是只作用于已经通过可行性校验的候选，不能覆盖工具不存在、schema
错误、I/O
不闭合和安全阻断；二是修正幅度存在上界，避免一次异常观测完全覆盖静态目标匹配和工程可靠性判断。借助这两个边界，算法在适应运行时变化的同时仍能保持可复现和可审计。

\subsection{恢复动作选择与 HITL
映射}\label{ux6062ux590dux52a8ux4f5cux9009ux62e9ux4e0e-hitl-ux6620ux5c04}

CEBRA-WP
将恢复动作限定为四类：\texttt{continue}、\texttt{patch\_local}、\texttt{suffix\_replan}
和 \texttt{stop}。四类动作与 FSM/HITL 的映射如表 4-3 所示。

\textbf{表 4-3 恢复动作与 FSM/HITL 映射}

{\def\LTcaptype{none} % do not increment counter
\begin{longtable}[]{@{}
  >{\raggedright\arraybackslash}p{(\linewidth - 8\tabcolsep) * \real{0.2000}}
  >{\raggedright\arraybackslash}p{(\linewidth - 8\tabcolsep) * \real{0.2000}}
  >{\raggedright\arraybackslash}p{(\linewidth - 8\tabcolsep) * \real{0.2000}}
  >{\raggedright\arraybackslash}p{(\linewidth - 8\tabcolsep) * \real{0.2000}}
  >{\raggedright\arraybackslash}p{(\linewidth - 8\tabcolsep) * \real{0.2000}}@{}}
\toprule\noalign{}
\begin{minipage}[b]{\linewidth}\raggedright
动作
\end{minipage} & \begin{minipage}[b]{\linewidth}\raggedright
触发背景
\end{minipage} & \begin{minipage}[b]{\linewidth}\raggedright
系统效果
\end{minipage} & \begin{minipage}[b]{\linewidth}\raggedright
FSM/HITL 映射
\end{minipage} & \begin{minipage}[b]{\linewidth}\raggedright
审计结果
\end{minipage} \\
\midrule\noalign{}
\endhead
\bottomrule\noalign{}
\endlastfoot
\texttt{continue} & 成功概率较高、证据充分、风险和预算压力可接受 &
继续执行当前计划后续步骤 & RUNNING 内自动推进或维持 PLANNED/RUNNING &
记录动作效用与继续原因 \\
\texttt{patch\_local} & 局部失败、重试耗尽、前缀仍有效、局部可修复性较高
& 生成局部 PlanPatch，替换参数或局部工具 & WAITING\_PATCH\_CONFIRM
或策略允许的受控自动 patch & 记录 patch 候选、失败上下文和应用结果 \\
\texttt{suffix\_replan} & 结构性失败概率升高、后缀可靠性不足、前缀可保留
& 保留已验证前缀，替换未执行后缀 & WAITING\_REPLAN\_CONFIRM & 记录
replan 候选、前缀保留位置和新后缀 \\
\texttt{stop} & 成功概率低、预算压力高、恢复余量低、人工介入价值有限 &
生成终止型 ReplanCandidate & WAITING\_REPLAN\_CONFIRM；接受后进入 FAILED
& 记录 terminal\_reason 与已保留证据 \\
\end{longtable}
}

动作效用由状态量和派生量共同计算。为便于表示，记
\texttt{s=p\_success}，\texttt{f=p\_structural\_failure}，\texttt{r=recovery\_margin}，\texttt{e=evidence\_sufficiency}，\texttt{b=budget\_pressure}。局部可修复性、证据可复用性、前缀可保留性、预算缓解度、目标重对齐收益和人工介入价值，分别记为
\texttt{local\_patchability}、\texttt{evidence\_reusability}、\texttt{prefix\_preservability}、\texttt{budget\_relief}、\texttt{goal\_realignment}
和 \texttt{intervention\_value}。动作效用可写为：

\begin{Shaded}
\begin{Highlighting}[]
\NormalTok{U\_continue}
\NormalTok{  = 0.38s + 0.14e + 0.12r {-} 0.22f {-} 0.14b}

\NormalTok{U\_patch\_local}
\NormalTok{  = 0.20s + 0.24r + 0.18 · local\_patchability}
\NormalTok{    + 0.12 · evidence\_reusability {-} 0.14f {-} 0.12b}

\NormalTok{U\_suffix\_replan}
\NormalTok{  = 0.18(1{-}s) + 0.20f + 0.16(1{-}r)}
\NormalTok{    + 0.18 · prefix\_preservability + 0.14 · budget\_relief}
\NormalTok{    + 0.14 · goal\_realignment}

\NormalTok{U\_stop}
\NormalTok{  = 0.32(1{-}s) + 0.24b + 0.18(1{-}r)}
\NormalTok{    + 0.16 · safety\_terminality + 0.10(1{-}intervention\_value)}
\end{Highlighting}
\end{Shaded}

这组效用函数体现了四类动作的不同偏好：\texttt{continue}
对较高成功概率、证据充分度和恢复余量更敏感；\texttt{patch\_local}
更依赖局部可修复性、证据可复用性和较低结构性失败概率；\texttt{suffix\_replan}
偏向于结构性失败概率较高、当前成功概率较低但前缀仍可保留的场景；\texttt{stop}
只有在继续价值低、预算压力高、恢复余量低且人工介入价值不足时才会成为强候选。

\texttt{stop}
不是用户主动取消，也不是执行异常崩溃，而是算法提出的终止型重规划候选。其载体为
\texttt{ReplanCandidate(replan\_mode="terminal\_stop")}，通过
\texttt{replan\_confirm} 通道进入
WAITING\_REPLAN\_CONFIRM。只有在人工接受或策略显式允许时，系统才将任务推进到
FAILED；已完成前缀、产物和解释字段仍被保留为审计资产。

综合上述设计，CEBRA-WP 主流程可写为算法 4-1。

\textbf{算法 4-1 CEBRA-WP 主流程}

\begin{Shaded}
\begin{Highlighting}[]
\NormalTok{Input:}
\NormalTok{  g: 设计目标}
\NormalTok{  C: 约束集合}
\NormalTok{  K: ProteinToolKG}
\NormalTok{  h\_t: 执行历史}
\NormalTok{  o\_t: 当前观测}
\NormalTok{  x\_t: Lite belief{-}state}

\NormalTok{Output:}
\NormalTok{  Decision\_t: 候选集合、默认候选、恢复动作、解释和证据引用}

\NormalTok{1.  Pi\_raw,t \textless{}{-} GenerateCandidates(g, C, K, h\_t)}
\NormalTok{2.  Pi\_t \textless{}{-} FeasibilityFilter(Pi\_raw,t, C, K, h\_t)}
\NormalTok{3.  if Pi\_t is empty and degraded\_feasible candidates exist:}
\NormalTok{4.      Pi\_t \textless{}{-} degraded\_feasible candidates}
\NormalTok{5.      mark candidates as requiring HITL confirmation}
\NormalTok{6.  for each pi in Pi\_t:}
\NormalTok{7.      S\_static(pi) \textless{}{-} StaticUtility(pi, g, C, K)}
\NormalTok{8.      S\_post(pi) \textless{}{-} ReplaceObjectiveTerm(S\_static(pi), PosteriorObjective(pi, g, o\_t))}
\NormalTok{9.  x\_t+1 \textless{}{-} BeliefUpdate(x\_t, o\_t, h\_t)}
\NormalTok{10. for each pi in Pi\_t:}
\NormalTok{11.     Delta(pi, x\_t+1) \textless{}{-} RuntimeAdjustment(pi, x\_t+1)}
\NormalTok{12.     U\_pi(pi, x\_t+1) \textless{}{-} clip(S\_post(pi) + Delta(pi, x\_t+1), 0, 1)}
\NormalTok{13. TopK\_t \textless{}{-} SelectDiverseTopK(Pi\_t, U\_pi, capability\_coverage)}
\NormalTok{14. U\_a \textless{}{-} ComputeActionUtility(continue, patch\_local, suffix\_replan, stop)}
\NormalTok{15. a\_t \textless{}{-} ApplyHardPrioritiesAndSelectAction(U\_a, x\_t+1, h\_t, C)}
\NormalTok{16. return Decision\_t(TopK\_t, a\_t, explanations, evidence\_refs)}
\end{Highlighting}
\end{Shaded}

算法第 1 至 8 行生成并评价候选，第 9 至 12
行根据运行时状态修正候选效用，第 13 行保留 Top-K 多样性，第 14 至 15
行选择恢复动作。第 15 行的硬优先级包括安全阻断优先、schema/I/O/tool
availability
违规淘汰、重试耗尽后优先局部修补、前缀可保留时优先后缀重规划等约束。借助这些优先级，动作选择遵循系统控制边界，而不是只按连续效用值排序。

\subsection{策略组与实验可验证性}\label{ux7b56ux7565ux7ec4ux4e0eux5b9eux9a8cux53efux9a8cux8bc1ux6027}

为验证 CEBRA-WP
各组成机制的作用，系统将算法能力拆分为四组策略开关。四组策略并非四套不同系统，而是在同一工程实现上逐步打开静态评分、固定阈值门控、动态恢复和
Lite belief-state 的内部消融配置，如表 4-4 所示。

\textbf{表 4-4 四组消融策略与算法机制开关}

{\def\LTcaptype{none} % do not increment counter
\begin{longtable}[]{@{}
  >{\raggedright\arraybackslash}p{(\linewidth - 14\tabcolsep) * \real{0.1250}}
  >{\raggedright\arraybackslash}p{(\linewidth - 14\tabcolsep) * \real{0.1250}}
  >{\raggedright\arraybackslash}p{(\linewidth - 14\tabcolsep) * \real{0.1250}}
  >{\raggedright\arraybackslash}p{(\linewidth - 14\tabcolsep) * \real{0.1250}}
  >{\raggedright\arraybackslash}p{(\linewidth - 14\tabcolsep) * \real{0.1250}}
  >{\raggedright\arraybackslash}p{(\linewidth - 14\tabcolsep) * \real{0.1250}}
  >{\raggedright\arraybackslash}p{(\linewidth - 14\tabcolsep) * \real{0.1250}}
  >{\raggedright\arraybackslash}p{(\linewidth - 14\tabcolsep) * \real{0.1250}}@{}}
\toprule\noalign{}
\begin{minipage}[b]{\linewidth}\raggedright
策略组
\end{minipage} & \begin{minipage}[b]{\linewidth}\raggedright
静态评分
\end{minipage} & \begin{minipage}[b]{\linewidth}\raggedright
固定阈值门控
\end{minipage} & \begin{minipage}[b]{\linewidth}\raggedright
动态观测
\end{minipage} & \begin{minipage}[b]{\linewidth}\raggedright
Lite belief-state
\end{minipage} & \begin{minipage}[b]{\linewidth}\raggedright
Runtime rerank
\end{minipage} & \begin{minipage}[b]{\linewidth}\raggedright
恢复动作效用
\end{minipage} & \begin{minipage}[b]{\linewidth}\raggedright
主要实验问题
\end{minipage} \\
\midrule\noalign{}
\endhead
\bottomrule\noalign{}
\endlastfoot
\texttt{static\_top1} & 启用 & 未启用 & 未启用 & 未启用 & 未启用 &
未启用 & 单一静态最优候选是否足以完成任务 \\
\texttt{fixed\_threshold\_gate} & 启用 & 启用 & 有限使用 & 未启用 &
未启用 & 未启用 & 固定阈值是否能控制风险与成本 \\
\texttt{dynamic\_no\_belief\_state} & 启用 & 启用 & 启用 & 未启用 &
部分启用 & 部分启用 & 直接运行时观测是否已能支撑恢复 \\
\texttt{lite\_belief\_state} & 启用 & 启用 & 启用 & 启用 & 启用 & 启用 &
显式状态是否带来更好的解释、预算感知和恢复决策 \\
\end{longtable}
}

该策略设计使第 7
章实验能够分别回答三个问题：静态计划是否足够；固定阈值门控是否会带来额外成本；在动态恢复已经存在时，Lite
belief-state 的增量价值体现在哪里。根据已有实验设置，本文不把 CEBRA-WP
论证为所有指标上的性能最优算法，而是重点验证其机制链路是否可执行、状态是否可观测、成本控制是否具有方向性证据，以及恢复决策是否具有可解释依据。

\subsection{数据契约与模块协作}\label{ux6570ux636eux5951ux7ea6ux4e0eux6a21ux5757ux534fux4f5c}

系统以统一数据契约支撑算法与工程实现之间的衔接。核心契约关系如图 4-7
所示。

图 4-7 核心数据契约 UML
插图文件：\texttt{paper/figures/uml-contracts.drawio.svg}

ProteinDesignTask 是任务入口契约，包含 task\_id、goal 和约束字段。Plan
由多个 PlanStep 组成，步骤之间通过 \texttt{S\{id\}.\{field\}}
引用语法建立数据依赖。StepResult
记录每个步骤的执行状态、输出、指标、产物路径和失败信息。PendingAction 与
Decision 构成 HITL
的双向契约：前者承载候选集合、默认建议和解释，后者记录用户选择、决策人和时间戳。TaskSnapshot
保存计划版本、已完成步骤、pending\_action\_id 和 runtime\_state
摘要。RuntimeState 持久化 CEBRA-WP 的五个核心状态量。DesignResult
是最终输出契约，包含序列、结构路径、评分、风险标记、报告路径和元数据。

图 4-8 给出了典型任务 t1 的贯穿示例。

图 4-8 t1 任务实例走查
插图文件：\texttt{paper/figures/t1-trpcage-instance-walkthrough.drawio.svg}

该示例从任务创建开始，经 Planner 生成候选计划、用户确认、Executor
调用工具、StepResult 写入、RuntimeState 更新和必要的恢复候选触发，最后由
Summarizer 汇总为 DesignResult。CEBRA-WP
在这一过程中并非孤立算法模块，而是嵌入任务生命周期的规划与恢复层：算法输出通过
PendingAction、PlanPatch、ReplanCandidate 和 RuntimeState
等契约进入系统，系统状态则仍由 FSM 和 Workflow 控制。

\subsection{本章小结}\label{ux672cux7ae0ux5c0fux7ed3}

本章完成了系统总体设计和 CEBRA-WP
算法定义。系统采用五层架构组织输入、规划、执行、安全汇总和资源管理，并通过
FSM 与 HITL
契约约束任务状态推进。蛋白质设计能力被组织为六阶段工作流，ProteinToolKG
则负责支持工具能力发现和 I/O 兼容校验。在算法层面，CEBRA-WP
将候选生成、硬可行性过滤、静态效用、后验目标评分、Lite
belief-state、运行时重排序和恢复动作选择组合成闭环，用于高代价、部分可观测、可失败的科研工作流恢复规划。下一章将在该设计基础上说明后端
API、工作流执行、RuntimeState 更新、工具适配和前端交互等工程实现。
